\section[Channel bounds with an explicit field resource]{Channel bounds: complete operators and a resource-dependent Maxwell theorem}
\label{ch:section}

An efficient channel is an input pattern whose selected output retains a
specified fraction of its incident power.  The question here is more
restrictive: every coherent combination of the counted inputs must also
carry positive power through the \emph{same} section, while the absolute
flux outside it stays small.  Counting individual basis fields is not enough,
because interference changes their flux.  Nor does a count of modes of an
empty aperture automatically control the fields inside an arbitrary device.

The route in this section makes explicit the missing physical resource.  A
bound on the field in a box around the section limits how many independent
inputs can each produce a prescribed through-flux.  The argument first
replaces the sharp section by a smooth window, then controls that window's
flux by the local magnetic field.  Maxwell's equations provide a derivative
bound on that field, allowing a finite-dimensional approximation in the box.
The resulting count depends on the local field resource; a resonator can make
that resource large even when its incident power is fixed.  Thus the estimate
is useful when the resource is separately controlled, but it is not an
area-only bound for the unrestricted class.

The unrestricted fixed-section assertion is disproved later in
Theorem~\ref{uc:counterexample} for $\delta>\eta$. The stronger small-bypass
regime of the resource bounds below remains a separate question.
This section proves three usable conclusions: an exact finite-rank reduction,
a local full-vector field-resource estimate, and a bridge from the prescribed
absolute-bypass condition to that estimate.  The local estimates reproduce,
with a complete audit, the argument in the supplied report
\cite[Section 3]{priorreport}.  The outer-window construction below avoids that
report's extra conserved-collar assumption for a sharp section.  None of these
results bounds resonant amplification by external passivity.

\subsection{What the cited channel bounds actually control}
\label{ch:literature}

Miller's Supplement S2 explicitly treats the maximum-index substitution as a
conjecture; S12 distinguishes homogeneous-background SVD estimates from a
design-specific calculation of the complete Green function.  These are useful
diagnostics, without a uniform theorem for every structured interior
\cite[Supplements S2 and S12]{miller2023}.

Kuang, Miller, and Miller first identify the physically normalized map as
$R_R^{1/2}T_RGT_SR_S^{-1/2}$.  Their Eq.~(1) retains the product of the two
endpoint operator norms when bounding its singular values by those of $G$.
Their explicit domain bounds concern the background Green operator, with
source/receiver size and separation data and spherical separation; their
relative-channel thresholds use a sum-rule normalization.  Appendix E rejects
the older largest-singular-value normalization as a general monotonic bound
\cite[Section II and Appendix E]{kuang2025}.  In particular, neither the
endpoint norms nor these geometric data are determined by the area of the
internal section.  A continuum use additionally needs bounded endpoint maps:
invertibility of a finite resistance matrix does not prove a uniformly bounded
inverse in an infinite-dimensional current space.

\subsection{Exact operator reduction}
\label{ch:operators}

We begin with a simple observation about the complete input-to-output map.
Suppose it can be approximated, uniformly on the candidate input space, by a
map with only $M$ output directions.  If there were more than $M$ independent
efficient inputs, some nonzero combination would disappear under the
approximating map.  The approximation error would then have to carry its
entire output, contradicting the efficiency threshold.  The following
lemma states this argument without assuming that the full map is compact.

Throughout, norms in the incident and delivered channel spaces are physical
power norms.  Let $T:\mathcal H_{\rm in}\to\mathcal H_{\rm out}$ be a bounded
operator, and fix $0<\tau<1$.  A counted linear subspace $W$ satisfies
\begin{equation}
 \|Ta\|^2\geq\tau\|a\|^2\qquad(a\in W).
 \label{ch:throughput}
\end{equation}
The other passage requirements are imposed separately below.

\begin{lemma}[PROVED: finite-rank obstruction]
\label{ch:ranklemma}
If a linear map $R:W\to\mathcal H_{\rm out}$ has rank at most $M$ and
$\|T|_W-R\|<\sqrt\tau$, then $\dim W\leq M$.
\end{lemma}
\begin{proof}
If $\dim W>M$, choose $M+1$ linearly independent vectors in $W$.
The restriction of $R$ to their span has a nonzero kernel vector $a$.
Then $\|Ta\|=\|(T-R)a\|<\sqrt\tau\|a\|$, contradicting
\eqref{ch:throughput}.  This proof also covers an infinite-dimensional $W$.
The strict inequality is indispensable: on $\mathbb C^{M+1}$,
$T=\sqrt\tau I$ and $R=0$ attain equality in the norm estimate.
\end{proof}

\begin{definition}[Approximation number]
\label{ch:approxdef}
For a bounded operator $U$ between Hilbert spaces, define
\begin{equation}
 a_{M+1}(U)=\inf_{\operatorname{rank}R\leq M}\|U-R\|.
 \label{ch:approxnumber}
\end{equation}
No compactness is required for this definition.  Thus
$a_{M+1}(T|_W)<\sqrt\tau$ is sufficient in Lemma~\ref{ch:ranklemma}.
\end{definition}

\begin{proposition}[PROVED: factorization with every amplification factor retained]
\label{ch:factorization}
Suppose $T=AGB+E$, where all maps are bounded between their stated Hilbert
spaces and $\|E\|\leq e$.  If $\operatorname{rank}G_M\leq M$, then
\begin{equation}
 \operatorname{rank}(AG_MB)\leq M,\qquad
 \|T-AG_MB\|\leq \|A\|\,\|B\|\,\|G-G_M\|+e.
 \label{ch:factorerror}
\end{equation}
Consequently $\dim W\leq M$ if the right side is less than $\sqrt\tau$.
If $G$ is compact, one may replace $\|G-G_M\|$ by its $(M+1)$st singular
value using the truncated singular-value expansion.
\end{proposition}
\begin{proof}
The image of $AG_MB$ is contained in the image under $A$ of the
at-most-$M$-dimensional image of $G_M$.  Submultiplicativity and the triangle
inequality prove \eqref{ch:factorerror}.  Apply Lemma~\ref{ch:ranklemma}.
For compact $G$, its singular expansion is an orthogonal expansion and the
norm of its discarded tail is exactly its next singular value.
\end{proof}

\begin{remark}[PROVED algebraic counteraudit; no Maxwell counterexample]
\label{ch:algebraaudit}
For $G_N=\operatorname{diag}(1,2^{-1},\ldots,N^{-1})$ and
$A_N=\operatorname{diag}(1,2,\ldots,N)$, the full map $A_NG_N=I_N$
is passive and has $N$ unit transmission singular values.  Passivity controls
the product, not $\|A_N\|=N$.  Likewise, the finite-rank orthogonal projections
onto the first $N$ coordinates of $\ell^2$ are individually compact
contractions, although their counts above any fixed threshold below one grow
without bound.  Finally, compressing $\operatorname{diag}(100,1,1)$ to the
last two coordinates decreases every singular value, but increases the ratio
of the second to the first from $1/100$ to $1$.  These exact examples identify
invalid operator inferences.  They supply no dielectric realization and no
section or bypass certificate.
\end{remark}

\subsection{Local Maxwell hypotheses and a legitimate weighted flux}
\label{ch:local}

To apply a finite-dimensional argument to passage through a section, we need
a measurement that is defined for the actual fields and can be estimated by
volume norms.  The smooth weight $\psi$ below measures flux through a window
on the section plane; it is a mathematical aid, not a change in the physical
square.  A cutoff in the longitudinal direction turns this weighted plane
measurement into a volume integral.  The term involving the transverse gradient
of $\psi$ records power crossing the side of that window, so conservation
does not require us to assume that this lateral flux vanishes.

Let $B\subset\mathbb R^3$ be an open rectangular box with side lengths
$L_1,L_2,L_3$ and volume $V_B$.  The physical dielectric extends beyond $B$;
no physical boundary condition is imposed on its faces.  A fixed well-posed
scattering problem supplies fields depending linearly on the input $a$.
Write
\begin{equation}
 e_a=\sqrt{\varepsilon_0}E_a,\qquad h_a=\sqrt{\mu_0}H_a,
 \qquad c_0=(\varepsilon_0\mu_0)^{-1/2},\qquad k_0=\omega/c_0.
 \label{ch:normalization}
\end{equation}
Assume on $B$, in distributions,
\begin{equation}
 \nabla\times e_a=ik_0h_a,\qquad
 \nabla\times h_a=-ik_0\varepsilon_r e_a,\qquad
 0<\varepsilon_-\leq\varepsilon_r(x)\leq\varepsilon_+<\infty.
 \label{ch:localmaxwell}
\end{equation}
The real scalar coefficient may be merely measurable.  Define the local
weighted field resource
\begin{equation}
 \mathcal E_B(a)=\int_B(\varepsilon_r|e_a|^2+|h_a|^2)\,dx.
 \label{ch:resource}
\end{equation}
Each field has finite resource.  A uniform bound relative to incident power
will be an \emph{additional} assumption.  For nondispersive material this
integral is four times the time-averaged local electromagnetic energy.
In dispersive material a physical stored-energy formula requires frequency
derivatives of the constitutive law, absent from these hypotheses.

Choose a real $\psi\in W^{1,\infty}_c(\mathbb R^2)$ with $0\leq\psi\leq1$,
and $\ell>0$, so that $B$ contains a neighborhood of
$\operatorname{supp}\psi\times[0,\ell]$, including a negative-$z$ collar.
Let $\chi\in W^{1,\infty}(0,\ell)$ be real, with
$\chi(0)=1$, $\chi(\ell)=0$.  Set
\begin{equation}
 g(x,y,z)=
 \begin{cases}
  (\chi(z)\nabla_\perp\psi(x,y),\chi'(z)\psi(x,y)),&0<z<\ell,\\
  0,&\text{otherwise},
 \end{cases}
 \qquad G=\|g\|_\infty.
 \label{ch:g}
\end{equation}
Select real cutoffs $0\leq\zeta,\rho\leq1$ in
$W^{1,\infty}_c(B)$, with $\zeta=1$ on $\operatorname{supp}g$ and
$\rho=1$ on $\operatorname{supp}\zeta$.  Introduce explicit constants
\begin{equation}
 \begin{split}
 d_1&=\|\nabla\rho\|_\infty,\qquad d_2=\|\nabla\zeta\|_\infty,\\
 C_H&=k_0^2\varepsilon_++d_1^2,\qquad
 C_E=\left(2+\frac{4d_2^2}{k_0^2\varepsilon_-}\right)^{1/2},\qquad
 K=\frac{c_0GC_E}{2\sqrt{\varepsilon_-}}.
 \end{split}
 \label{ch:constants}
\end{equation}

\begin{lemma}[PROVED: canonical weak flux]
\label{ch:fluxlemma}
Define the Hermitian cross-Poynting vector
\begin{equation}
 p(a,b)=\frac{c_0}{4}
 \left(e_a\times\overline{h_b}+\overline{e_b}\times h_a\right).
 \label{ch:crosspoynting}
\end{equation}
For almost every $z$ near $[0,\ell]$, put
$F_{ab}(z)=\int_{\mathbb R^2}\psi p_z(a,b)\,dx\,dy$.
This is a $W^{1,1}$ function of $z$, with its unique continuous
representative satisfying
\begin{equation}
 F'_{ab}(z)=\int_{\mathbb R^2}\nabla_\perp\psi\cdot p_\perp(a,b)\,dx\,dy,
 \qquad
 q_\psi(a,b):=F_{ab}(0)=-\int_Bg\cdot p(a,b)\,dx.
 \label{ch:weakflux}
\end{equation}
The diagonal $p(a,a)$ is the physical real time-averaged Poynting vector.
\end{lemma}
\begin{proof}
Both fields are locally in $H(\operatorname{curl})$ by
\eqref{ch:localmaxwell}; their cross products belong to $L^1$.
For such fields the identity
$\operatorname{div}(u\times v)=v\cdot\operatorname{curl}u-
u\cdot\operatorname{curl}v$ holds in distributions.  To justify it, mollify
both fields in a relatively compact neighborhood of the test function;
the fields and their curls converge in $L^2$, and every product converges
in $L^1$.  For the first cross product in \eqref{ch:crosspoynting} its
divergence is
$ik_0(h_a\cdot\overline{h_b}-\varepsilon_r e_a\cdot\overline{e_b})$.
The second has the opposite divergence, since $\varepsilon_r$ is real.
Thus $\operatorname{div}p(a,b)=0$.  Testing this with
$\psi(x,y)\theta(z)$ and using Fubini proves the derivative formula in
\eqref{ch:weakflux}.  Both terms are integrable in $z$, establishing
$W^{1,1}$ membership.  Finally integrate $(\chi F_{ab})'$ between zero
and $\ell$ and use the endpoint values of $\chi$.  This gives the volume
formula, including its transverse-gradient contribution.
\end{proof}

The transverse-gradient contribution accounts for lateral collar flux.
Dropping it would assert an unproved conservation law for the individual
aperture.  This construction requires no $L^2$ tangential trace of a rough
Maxwell field.

\subsection{Magnetic compactness and its quantitative cost}
\label{ch:magnetic}

The counting step uses the magnetic field because its curl and divergence
give a local first-derivative estimate even when the dielectric coefficient
is only measurable.  The cutoff $\rho$ localizes the field inside the box.
We first bound its gradient by the local field resource, then show that its
$L^2$ norm controls the weighted flux.  Together, these estimates prevent an
input with substantial passage and bounded resource from living entirely in
arbitrarily fine box eigenfunctions.

\begin{lemma}[PROVED: uniform magnetic regularity]
\label{ch:magneticHone}
The zero extension of $\rho h_a$ is in $H^1_0(B)^3$, and
\begin{equation}
 \|\nabla(\rho h_a)\|_2^2\leq C_H\mathcal E_B(a).
 \label{ch:Honebound}
\end{equation}
\end{lemma}
\begin{proof}
Taking divergence of $\operatorname{curl}e_a=ik_0h_a$ gives
$\operatorname{div}h_a=0$.  The compactly supported zero extension
$v=\rho h_a$ has curl and divergence in $L^2$.  The Fourier identity
$|\xi|^2|\widehat v|^2=|\xi\times\widehat v|^2+
|\xi\cdot\widehat v|^2$ gives
\begin{equation}
 \|\nabla v\|_2^2=\|\operatorname{curl}v\|_2^2+
 \|\operatorname{div}v\|_2^2.
 \label{ch:divcurl}
\end{equation}
It proves $H^1$ membership; compact support supplies zero boundary trace.
Let $X^2=\int_B\varepsilon_r|e_a|^2$ and $Y^2=\int_B|h_a|^2$.
Minkowski's inequality in the direct sum of the curl and divergence spaces,
and $|u\times v|^2+|u\cdot v|^2=|u|^2|v|^2$ for real $u$, give
\begin{equation}
 \|\nabla(\rho h_a)\|_2
 \leq k_0\sqrt{\varepsilon_+}X+d_1Y
 \leq\sqrt{C_H}\sqrt{X^2+Y^2}.
 \label{ch:Honeproof}
\end{equation}
No derivative of $\varepsilon_r$ occurs.
\end{proof}

\begin{lemma}[PROVED: Caccioppoli and flux domination]
\label{ch:caccioppoli}
The complete physical fields obey
\begin{equation}
 \int_{\operatorname{supp}g}\varepsilon_r|e_a|^2
 \leq C_E^2\int_{\operatorname{supp}\zeta}|h_a|^2,
 \qquad
 |q_\psi(a,a)|\leq K\|\rho h_a\|_2^2.
 \label{ch:fluxdomination}
\end{equation}
\end{lemma}
\begin{proof}
Eliminating $e_a$ gives the distributional equation
\begin{equation}
 \operatorname{curl}(\varepsilon_r^{-1}\operatorname{curl}h_a)
 =k_0^2h_a.
 \label{ch:magneticsecond}
\end{equation}
Test with $\zeta^2h_a$, permissible by compactly supported
$H(\operatorname{curl})$ approximation.  Write
$A_0=\int_B\zeta^2\varepsilon_r^{-1}|\operatorname{curl}h_a|^2$ and
$J_0=\int_B\varepsilon_r^{-1}|\nabla\zeta\times h_a|^2$.
Taking real parts and applying Young's inequality gives
\begin{equation}
 A_0\leq k_0^2\int_B\zeta^2|h_a|^2+2\sqrt{A_0J_0}
 \leq k_0^2\int_B\zeta^2|h_a|^2+\tfrac12 A_0+2J_0.
 \label{ch:caccioppoliproof}
\end{equation}
Use $\varepsilon_r^{-1}\leq\varepsilon_-^{-1}$ and
$\varepsilon_r|e_a|^2=k_0^{-2}\varepsilon_r^{-1}
|\operatorname{curl}h_a|^2$ to obtain the first estimate in
\eqref{ch:fluxdomination}.  The volume representation and
Cauchy--Schwarz then give
\begin{equation}
 |q_\psi(a,a)|\leq\frac{c_0G}{2}
 \|e_a\|_{L^2(\operatorname{supp}g)}
 \|h_a\|_{L^2(\operatorname{supp}g)}
 \leq\frac{c_0GC_E}{2\sqrt{\varepsilon_-}}
 \|h_a\|_{L^2(\operatorname{supp}\zeta)}^2
 \leq K\|\rho h_a\|_2^2.
 \label{ch:fluxproof}
\end{equation}
This is a vector Maxwell estimate, including the actual reflected and
evanescent fields, not a scalar reduction or a propagating-wave cutoff.
\end{proof}

The next theorem implements that idea.  If the input space had more
dimensions than the retained box eigenfunctions, a nonzero combination would
have no component in any of those low-frequency directions.  Its gradient
bound would make its remaining magnetic field too small to carry the required
flux.  The count is therefore finite at fixed $B_W$, with an explicit cost
when $B_W$ increases.

\begin{theorem}[PROVED: through-flux resource bound]
\label{ch:resourcecount}
Let $W$ have finite dimension $N$, and suppose
$q_\psi(a,a)\geq p\|a\|^2$ for every $a\in W$, where $p>0$.
Define
\begin{equation}
 B_W=\sup_{a\in W,\,\|a\|=1}\mathcal E_B(a),\qquad
 \Lambda=K C_H B_W/p.
 \label{ch:cutoff}
\end{equation}
Then, with $N_{D,B}(\Lambda)$ the scalar Dirichlet eigenvalue count of the
comparison box,
\begin{equation}
 N\leq 3N_{D,B}(\Lambda)
 \leq\frac{V_B}{2\pi^2}\left(\frac{K C_H B_W}{p}\right)^{3/2}.
 \label{ch:count}
\end{equation}
Here $B_W$ is finite for a finite-dimensional $W$, but it is not uniformly
bounded by the fixed physical class.
\end{theorem}
\begin{proof}
Let $\Pi_\Lambda$ project onto vector Dirichlet eigenfunctions with scalar
eigenvalue at most $\Lambda$.  Its rank is $3N_{D,B}(\Lambda)$.
If $N$ exceeds that rank, there is a unit $a\in W$ with
$\Pi_\Lambda(\rho h_a)=0$.  Writing $\lambda_{\rm next}>\Lambda$ for
the next scalar eigenvalue, Parseval and \eqref{ch:Honebound} imply
\begin{equation}
 \|\rho h_a\|_2^2\leq
 \frac{\|\nabla(\rho h_a)\|_2^2}{\lambda_{\rm next}}
 \leq\frac{C_HB_W}{\lambda_{\rm next}}.
 \label{ch:tail}
\end{equation}
Equation~\eqref{ch:fluxdomination} contradicts $q_\psi(a,a)\geq p$.
The box eigenvalues are
\begin{equation}
 \lambda_{\mathbf n}=\pi^2\sum_{i=1}^3 n_i^2/L_i^2,
 \qquad\mathbf n\in\{1,2,\ldots\}^3.
 \label{ch:boxspectrum}
\end{equation}
For each point below the cutoff, the unit cell
$\prod_i[n_i-1,n_i]$ lies in the positive octant of
$\sum_i x_i^2/L_i^2\leq\Lambda/\pi^2$.  These cells have disjoint
interiors.  Their total volume is at most the ellipsoid-octant volume,
giving $N_{D,B}(\Lambda)\leq V_B\Lambda^{3/2}/(6\pi^2)$.
Multiplication by three proves \eqref{ch:count} without a Weyl asymptotic
or a remainder depending on fine material interfaces.
\end{proof}

\begin{theorem}[PROVED: cost averaged over an orthonormal input basis]
\label{ch:average}
Under the hypotheses of Theorem~\ref{ch:resourcecount}, let
$a_1,\ldots,a_N$ be any orthonormal basis of $W$, and set
$\mathcal E_{\rm sum}=\sum_i\mathcal E_B(a_i)$,
$B_{\rm avg}=\mathcal E_{\rm sum}/N$.  If $\lambda_j^{\rm vec}$ lists
the Dirichlet spectrum with three component multiplicity, then
\begin{equation}
 \mathcal E_{\rm sum}\geq\frac{p}{K C_H}
 \sum_{j=1}^N\lambda_j^{\rm vec}
 \geq\frac{3p}{5K C_H}
 \left(\frac{2\pi^2}{V_B}\right)^{2/3}N^{5/3},\qquad
 N\leq\frac{V_B}{2\pi^2}
 \left(\frac{5K C_HB_{\rm avg}}{3p}\right)^{3/2}.
 \label{ch:averagebound}
\end{equation}
\end{theorem}
\begin{proof}
Let $A:\mathbb C^N\to H^1_0(B)^3$ take coefficients to
$\rho h_{\sum_i c_i a_i}$.  Its $L^2$ Gram matrix $S=A^*A$ obeys
$S\geq(p/K)I$, by flux coercivity and \eqref{ch:fluxdomination}.
The columns $u_i$ of $AS^{-1/2}$ are $L^2$-orthonormal.  Let $D_\nabla$
be their positive semidefinite gradient Gram matrix.  Cyclicity of finite
matrix trace gives
\begin{equation}
 \sum_i\|\nabla(\rho h_{a_i})\|_2^2
 =\operatorname{Tr}(S D_\nabla)
 \geq\frac pK\operatorname{Tr}D_\nabla
 \geq\frac pK\sum_{j=1}^N\lambda_j^{\rm vec}.
 \label{ch:averageproof}
\end{equation}
For the last inequality, expand the $u_i$ in a vector Dirichlet eigenbasis
$v_j$.  The weights $w_j=\sum_i|\langle u_i,v_j\rangle|^2$ are between
zero and one by Bessel's inequality and sum to $N$.  Among such weights,
$\sum_j\lambda_j^{\rm vec}w_j$ is minimized by filling the first $N$
positions; equivalently compare all weights with the cutoff
$\lambda_N^{\rm vec}$.  All series are nonnegative.
The left side of \eqref{ch:averageproof} is at most
$C_H\mathcal E_{\rm sum}$.  The box count gives
$\lambda_j^{\rm vec}\geq(2\pi^2j/V_B)^{2/3}$, and
$\sum_{j=1}^Nj^{2/3}\geq\int_0^Nx^{2/3}dx=3N^{5/3}/5$.
This proves both inequalities.  Basis independence of
$\mathcal E_{\rm sum}$ follows because it is the trace of the positive
resource quadratic form restricted to $W$.
\end{proof}

\subsection{Absolute bypass supplies the sharp-section bridge}
\label{ch:bypassbridge}

The preceding theorem concerns a smooth window, whereas the question is
posed for the sharp square $D$.  Choose a window equal to one on $D$ and
tapering to zero outside it.  The correction introduced by the taper can
have either sign, but its magnitude is at most the allowed absolute bypass.
Passage and throughput supply the positive contribution $\eta\tau$ per unit
incident power, leaving a positive margin precisely when the sufficient
condition $\delta<\eta\tau$ holds.  This is a condition for this bridge;
it does not settle the other threshold regimes.

The following trace hypothesis states explicitly how the physical section
measurement is related to Lemma~\ref{ch:fluxlemma}.  It is automatic for
smooth fields on a neighborhood of the section, and also whenever the
relevant plane Poynting integrals converge in $L^1$ to their plane-$0$
values.  It is not inferred from $H(\operatorname{curl})$ alone.

\begin{assumption}[Physical section compatibility]
\label{ch:tracecompatibility}
For each $a\in W$ there is a real flux density $s_a$ on the entire section
plane, integrable on $D$ and absolutely integrable on its complement,
such that
\begin{equation}
 P_D(a)=\int_Ds_a,\qquad
 q_\psi(a,a)=\int_{\mathbb R^2}\psi s_a,\qquad
 \int_{\mathbb R^2\setminus D}|s_a|\leq\delta\|a\|^2.
 \label{ch:traceagreement}
\end{equation}
Here $s_a=\tfrac12\operatorname{Re}(E_a\times\overline{H_a})\cdot\hat z$
when that classical trace exists.  A separately specified regularized
measurement may be used if it establishes these identities.
\end{assumption}

\begin{lemma}[PROVED: an outer window preserves passage up to bypass]
\label{ch:outerwindow}
Take $\psi=1$ on the same fixed section $D$, with $0\leq\psi\leq1$
and compact support in a slightly larger transverse region.  Under
Assumption~\ref{ch:tracecompatibility} and the physical passage condition
$P_D(a)\geq\eta\|Ta\|^2$, one has, for every $a\in W$,
\begin{equation}
 q_\psi(a,a)\geq\eta\|Ta\|^2-\delta\|a\|^2.
 \label{ch:outerwindowbound}
\end{equation}
If also \eqref{ch:throughput} holds and $\delta<\eta\tau$, then
$q_\psi(a,a)\geq p\|a\|^2$ with $p=\eta\tau-\delta>0$.
\end{lemma}
\begin{proof}
Because $\psi=1$ on $D$,
$q_\psi(a,a)=P_D(a)+\int_{\mathbb R^2\setminus D}\psi s_a$.
The second integral is at least
$-\int_{\mathbb R^2\setminus D}|s_a|$.
Apply the passage and bypass hypotheses, then the throughput bound.
\end{proof}

\begin{remark}[Scope of the outer-window argument]
\label{ch:outerwindowscope}
The physical section $D$ has not been enlarged: its prescribed bypass
condition controls the correction outside it.  The window is only an
analytical device.  No limit of a discontinuous indicator, no unproved
zero lateral flux, and no assumption about how much flux avoids an inner
edge strip are used.  The new restriction $\delta<\eta\tau$ is genuine
for this coercivity argument.  For larger tolerances this proof yields no
positive lower bound, and makes no claim of impossibility or of a
counterexample.  A sharp flux defined on an exceptional plane without the
compatibility in Assumption~\ref{ch:tracecompatibility} is likewise outside
the conclusion.
\end{remark}

\begin{theorem}[PROVED under stated resources: approximation of the complete transmission map]
\label{ch:fulloutput}
Let $W$ be a closed input subspace satisfying passage and absolute bypass,
Assumption~\ref{ch:tracecompatibility}, and the preceding local Maxwell
hypotheses.  Suppose for some finite $B_*$,
\begin{equation}
 \mathcal E_B(a)\leq B_*\|a\|^2\qquad(a\in W).
 \label{ch:uniformresource}
\end{equation}
For any $\Lambda\geq0$, let
$L_\Lambda a=\Pi_\Lambda(\rho h_a)$, let
$\mathcal K_\Lambda=\ker L_\Lambda$, and let $P_\Lambda$ be the
orthogonal projection of $W$ onto $\mathcal K_\Lambda^\perp$.
The operator $R_\Lambda=T|_W P_\Lambda$ satisfies
\begin{equation}
 \operatorname{rank}R_\Lambda\leq3N_{D,B}(\Lambda),\qquad
 \|T|_W-R_\Lambda\|^2
 \leq\frac1\eta\left(\delta+
             \frac{K C_H B_*}{\lambda_{\rm next}}\right),
 \label{ch:fulloutputerror}
\end{equation}
where $\lambda_{\rm next}$ is the first scalar box eigenvalue exceeding
$\Lambda$.  In particular, for $p=\eta\tau-\delta>0$ and
$\Lambda=K C_HB_*/p$, the norm error is strictly below $\sqrt\tau$.
If every vector in $W$ also satisfies \eqref{ch:throughput}, then
\begin{equation}
 \dim W\leq\frac{V_B}{2\pi^2}
       \left(\frac{K C_HB_*}{\eta\tau-\delta}\right)^{3/2}.
 \label{ch:physicalcount}
\end{equation}
\end{theorem}
\begin{proof}
The resource bound makes $a\mapsto\rho h_a$ bounded, so
$\mathcal K_\Lambda$ is closed.  Since $L_\Lambda$ has rank at most
$3N_{D,B}(\Lambda)$, its kernel has codimension at most that number;
this proves the rank statement.  For $b\in\mathcal K_\Lambda$,
the spectral-tail estimate and Lemmas~\ref{ch:caccioppoli}
and \ref{ch:outerwindow} give
\begin{equation}
 \eta\|Tb\|^2\leq q_\psi(b,b)+\delta\|b\|^2
 \leq\left(\frac{K C_HB_*}{\lambda_{\rm next}}+\delta\right)\|b\|^2.
 \label{ch:kerneloutput}
\end{equation}
Apply this to $b=(I-P_\Lambda)a$, whose norm is at most $\|a\|$.
Because $\lambda_{\rm next}>\Lambda=K C_HB_*/p$, the error is
strictly below $\sqrt\tau$ (if $B_*=0$, positivity of $p$ already
precludes a nonzero counted subspace).  Lemma~\ref{ch:ranklemma} and
the explicit box count finish the proof.
\end{proof}

\begin{corollary}[CONDITIONAL as a uniform family bound]
\label{ch:spectralcor}
Let $V_\tau=\operatorname{Ran}\mathbf1_{[\tau,\infty)}(T^*T)$.
If passage, absolute bypass, trace compatibility, common local geometry,
and a common finite resource $B_*$ hold for every vector of $V_\tau$,
with $\delta<\eta\tau$, then \eqref{ch:physicalcount} bounds
$\dim V_\tau$, without an a priori compactness assumption on $T$.
It is the usual count of squared singular values at least $\tau$ when
$T$ is compact.  Alternatively a uniform bound on the basis-average
resource over every finite-dimensional subspace gives the last bound in
\eqref{ch:averagebound}, with $p=\eta\tau-\delta$.
\end{corollary}
\begin{proof}
The spectral theorem implies \eqref{ch:throughput} on $V_\tau$.
Apply Theorem~\ref{ch:fulloutput}, or apply Theorem~\ref{ch:average}
to every finite-dimensional subspace.  An infinite-dimensional Hilbert
space contains finite orthonormal sets of arbitrarily large cardinality,
so a common finite upper bound forces finite dimension.
\end{proof}

The map $R_\Lambda$ is defined using the complete $T$, including all
matching and resonant feedback.  It is not obtained by deleting a component
of an isolated link and then silently applying an uncontrolled inverse.
Its cost is the explicitly retained norm $B_*$ of the actual full field.
Thus Theorem~\ref{ch:fulloutput} discharges the output-approximation step
under a resource hypothesis; it does not discharge that hypothesis.

\begin{corollary}[PROVED: essential-norm consequence]
\label{ch:essential}
Under Theorem~\ref{ch:fulloutput}'s hypotheses, without assuming a
throughput lower bound, the distance of $T|_W$ from the compact operators
is at most $\sqrt{\delta/\eta}$.  In particular, exact zero bypass makes
$T|_W$ compact.
\end{corollary}
\begin{proof}
The finite-rank $R_\Lambda$ are compact and
$\lambda_{\rm next}\to\infty$ as $\Lambda\to\infty$.
Take the limit superior in \eqref{ch:fulloutputerror}.  When $\delta=0$,
these finite-rank operators converge to $T|_W$ in operator norm.
A norm limit of compact operators is compact: given any bounded input
sequence, use a diagonal subsequence convergent under each approximant,
then the uniform approximation makes its images under $T$ Cauchy.
\end{proof}

\subsection{An explicit area scaling with its extra resource exposed}
\label{ch:explicitconstants}

\begin{corollary}[CONDITIONAL: square-enclosure area scaling]
\label{ch:arearesource}
Suppose $D\subset[-L/2,L/2]^2$ and there is a source-free comparison box
\begin{equation}
 B=(-L/2-4r,L/2+4r)^2\times(-3r,4r),\qquad r>0,
 \label{ch:explicitbox}
\end{equation}
on which the material bounds and \eqref{ch:uniformresource} hold.
There are admissible cutoffs with
\begin{equation}
 V_B=7r(L+8r)^2,\quad
 G\leq\sqrt3/r,\quad d_1,d_2\leq\sqrt3/r,\quad
 C_H\leq k_0^2\varepsilon_++3/r^2,\quad
 K\leq\frac{c_0\sqrt3}{2r\sqrt{\varepsilon_-}}
       \sqrt{2+\frac{12}{k_0^2\varepsilon_-r^2}}.
 \label{ch:explicitconstantsvalue}
\end{equation}
In particular, if $r=(k_0\sqrt{\varepsilon_+})^{-1}$ is admissible and
$\kappa=\varepsilon_+/\varepsilon_-$, then, with
$p=\eta\tau-\delta>0$,
\begin{equation}
 \dim W\leq\frac{7}{2\pi^2}
 (k_0\sqrt{\varepsilon_+}L+8)^2
 \left(2\sqrt3\sqrt\kappa\sqrt{2+12\kappa}\,
       \frac{\omega B_*}{p}\right)^{3/2}.
 \label{ch:arearesourcebound}
\end{equation}
For a scaled fixed shape with $L^2\leq c_D|D|$, this is
$O(1+k_0^2\varepsilon_+|D|)$ \emph{only if} $\omega B_*$, the material
ratio, thresholds, shape constant, and the required collar remain controlled.
\end{corollary}
\begin{proof}
For a one-dimensional interval use the function equal to one on that
interval, decreasing linearly to zero across an exterior strip of width
$r$, and zero beyond the strip.  Its Lipschitz constant is $1/r$.
Let $Q_s=[-L/2-s,L/2+s]^2$.  Take $\psi$ as the product of the two
one-dimensional cutoffs equal to one on $Q_0$ and zero outside $Q_r$,
and take $\chi(z)=1-z/r$ on $(0,r)$.  Then
$|\nabla_\perp\psi|\leq\sqrt2/r$ and $|g|\leq\sqrt3/r$.
Take $\zeta$ as the product cutoff equal to one on
$Q_r\times[0,r]$, supported on $Q_{2r}\times[-r,2r]$;
take $\rho$ equal to one on the latter box and supported on
$Q_{3r}\times[-2r,3r]$.  Each has gradient bounded by $\sqrt3/r$,
and both supports are compact in \eqref{ch:explicitbox}.
Equality on the indicated closed supports suffices in the preceding
estimates.  This proves \eqref{ch:explicitconstantsvalue}.
Substitute $r=(k_0\sqrt{\varepsilon_+})^{-1}$ in
\eqref{ch:physicalcount}, using $\omega=c_0k_0$; straightforward
cancellation gives \eqref{ch:arearesourcebound}.
Finally $(x+8)^2\leq2x^2+128$ proves the asserted order with every
additional dependency retained.
\end{proof}

\begin{corollary}[PROVED: local extension to real symmetric anisotropy]
\label{ch:anisotropy}
All local estimates and their conditional channel consequences remain
valid for measurable real symmetric tensors $\varepsilon_r(x)$ satisfying
$\varepsilon_-I\leq\varepsilon_r\leq\varepsilon_+I$, provided
$\varepsilon_r|e|^2$ is replaced by
$\overline e\cdot\varepsilon_r e$ in the resource.
\end{corollary}
\begin{proof}
Symmetry and reality give the same cross-Poynting cancellation.
The magnetic estimate uses
$|\varepsilon_r e|^2\leq\varepsilon_+
\overline e\cdot\varepsilon_r e$.
In the Caccioppoli step use the positive quadratic form of
$\varepsilon_r^{-1}$, its Cauchy--Schwarz inequality, and
$\varepsilon_r^{-1}\leq\varepsilon_-^{-1}I$.
The identity
$k_0^{-2}\overline{\operatorname{curl}h}\cdot
\varepsilon_r^{-1}\operatorname{curl}h=
\overline e\cdot\varepsilon_r e$ replaces its scalar version.
Every remaining step is unchanged.  Constant scalar permeability is still
essential to this proof of $\operatorname{div}h=0$.
\end{proof}

\subsection{The scope of the resource route}
\label{ch:open}

\begin{remark}[OPEN: removing the resource in the small-bypass regime]
\label{ch:missing}
There is no estimate here, from $1\leq\varepsilon_r\leq\varepsilon_+$,
the fixed section, passivity, passage, and absolute bypass alone, that
uniformly bounds $B_*$ or replaces its role by a uniform approximation
estimate for the full $T$.  A possible sufficient missing theorem is:
every admitted counted subspace has a finite-rank approximation with error
strictly below $\sqrt\tau$ and rank bounded solely by the frozen parameters.
An area scaling additionally requires a rank of order
$1+k_0^2\varepsilon_+|D|$.  Theorem~\ref{ch:fulloutput} supplies such an
approximation \emph{with} $B_*$; it does not eliminate it.  Even proving
unbounded local amplification would only invalidate that proof route,
not construct an arbitrarily large channel subspace.  Conversely,
pointwise compactness for each device would not make the bound uniform
over a sequence of devices. Theorem~\ref{uc:counterexample} later disproves
the unrestricted assertion in the allowed range $\delta>\eta$; it leaves
the stricter regime $\delta<\eta\tau$ used here unresolved without resources.
\end{remark}
